\begin{exercise}[Polymorphisme du carbonate de calcium]

Les deux formes allotropiques de \ce{CaCO3(s)} sont : la calcite et l'aragonite. Dans les conditions standards, les entropies molaires absolues de la calcite et de l'aragonite sont respectivement : \SI{92.80}{\joule\per\mole\per\kelvin} et \SI{88.62}{\joule\per\mole\per\kelvin}.


Leurs enthalpies molaires standards de formation sont respectivement : \SI{-1205.72}{\kilo\joule\per\mole} et \SI{-1205.90}{\kilo\joule\per\mole}. La transition de la calcite à l'aragonite se fait avec diminution de volume de \SI{2.75}{\centi\meter\cubed\per\mole}.

\begin{questions}
    \question Déterminer l'accroissement d'enthalpie libre pour la transition calcite $\to$ aragonite à \SI{25}{\celsius} sous une pression d'une atmosphère.
    
    \question Laquelle des deux formes est la plus stable dans ces conditions ?
    
    \question De combien faut-il accroître la pression, la température restant constante, pour que l'autre forme devienne stable ?
\end{questions}
\end{exercise}